\chapter[2000 --- Carlsson et al.]{2000: Arvid Carlsson, Paul Greengard, Eric R. Kandel}
\label{ch:2000_Carlsson}

\section*{Main Contribution}

\textit{Elucidated signal transduction in the nervous system, explaining how neurotransmitters trigger cascades of molecular events that strengthen or weaken synaptic connections.}

\section{Official Citation}

for their work concerning signal transduction in the nervous system

\section{Background and Historical Context}

The nervous system communicates through an intricate dance of electrical and chemical signals. Neurons transmit information as electrical impulses along their axons, and then convert these impulses into chemical signals---neurotransmitters---that cross the synaptic cleft and bind to receptors on postsynaptic neurons. The mechanisms by which neurotransmitters produce their effects inside target cells---the ``signal transduction'' pathways---were largely unknown until the work of Arvid Carlsson, Paul Greengard, and Eric R. Kandel, whose discoveries in the second half of the twentieth century earned them the 2000 Nobel Prize in Physiology or Medicine.

Arvid Carlsson was born on January 25, 1923, in Uddevalla, Sweden. He studied medicine at the University of Lund and earned his Ph.D.\ in 1959. Carlsson became professor of pharmacology at the University of Gothenburg, where he spent the majority of his career. His early research concerned the biochemistry of catecholamines---dopamine, norepinephrine, and epinephrine---and their role in the brain.

Paul Greengard was born on December 11, 1922, in New York City. He studied physics at Hamilton College and then earned his Ph.D.\ in biophysics at the Johns Hopkins University in 1953. After postdoctoral work at the National Institutes of Health and at the Medical Research Council in Cambridge, England, Greengard joined the faculty at Yale University and later at Rockefeller University, where he conducted his Nobel Prize--winning research on the intracellular signaling cascades activated by neurotransmitters.

Eric Richard Kandel was born on November 7, 1929, in Vienna, Austria. His family fled to the United States in 1939 to escape the Nazis. Kandel studied medicine at New York University, earning his M.D.\ in 1956, and then trained in psychiatry at the Massachusetts Mental Health Center and at the National Institute of Mental Health. He joined the faculty at Columbia University in 1974, where he became professor of biochemistry and biophysics and director of the Kavli Institute for Brain Science.

The intellectual context of their work was the revolution in molecular neuroscience that began in the 1950s and 1960s. The identification of neurotransmitters---acetylcholine, dopamine, norepinephrine, serotonin, and others---and the development of methods for studying synaptic transmission at the molecular level provided the tools for understanding how chemical signals are converted into cellular responses. However, the question of how neurotransmitter binding to receptors on the cell surface produced specific effects inside the cell---changes in ion channel activity, gene expression, and synaptic strength---remained to be answered.

\section{The Discovery/Contribution}

\subsection{Arvid Carlsson and Dopamine as a Neurotransmitter}

Carlsson's first major contribution was the demonstration that dopamine---a catecholamine previously thought to function only as a precursor to norepinephrine---was itself a neurotransmitter in the brain. In the late 1950s, Carlsson and his colleagues showed that dopamine was present in high concentrations in the basal ganglia of the brain and that it was released in response to nerve stimulation. They developed methods for measuring dopamine levels in brain tissue using fluorometric assays and demonstrated that dopamine-containing neurons formed a distinct neurotransmitter system---the nigrostriatal pathway---that was essential for the control of movement.

The clinical significance of this discovery became apparent when Carlsson linked dopamine dysfunction to Parkinson's disease, a debilitating movement disorder characterized by tremor, rigidity, and bradykinesia. Carlsson showed that the loss of dopamine-producing neurons in the substantia nigra was the primary pathological feature of Parkinson's disease, and he proposed that replacing dopamine could alleviate symptoms. This led to the development of L-DOPA (levodopa), a dopamine precursor that crosses the blood--brain barrier and is converted to dopamine in the brain. L-DOPA, first used clinically in the late 1960s, remains the most effective treatment for Parkinson's disease.

Carlsson also linked excess dopamine activity to schizophrenia and proposed the ``dopamine hypothesis'' of schizophrenia, which holds that the positive symptoms of the disorder---hallucinations and delusions---result from overactive dopamine signaling in the mesolimbic pathway. This hypothesis guided the development of antipsychotic drugs, which block dopamine D2 receptors in the brain.

\subsection{Paul Greengard and the Second Messenger Cascades}

Greengard's contribution was the elucidation of the intracellular signaling cascades activated by neurotransmitters. He showed that neurotransmitters bind to two major classes of receptors on the postsynaptic membrane: ionotropic receptors, which are ligand-gated ion channels that produce fast synaptic responses, and metabotropic receptors, which are G protein--coupled receptors that activate intracellular second messenger cascades producing slower, longer-lasting effects.

Greengard demonstrated that the neurotransmitter acetylcholine, acting through muscarinic (metabotropic) receptors, and other neurotransmitters acting through similar receptors, stimulate the production of second messengers---cyclic AMP (cAMP), cyclic GMP (cGMP), and inositol trisphosphate (IP3)---inside the postsynaptic cell. These second messengers activate specific protein kinases---protein kinase A (PKA), protein kinase G (PKG), and protein kinase C (PKC)---which phosphorylate ion channels and other target proteins, producing changes in the electrical properties of the neuron.

Greengard's primary discovery concerned DARPP-32 (dopamine- and cAMP-regulated phosphoprotein of molecular weight 32 kDa), a protein he identified in dopamine-responsive neurons. DARPP-32 is phosphorylated by PKA when dopamine is present and dephosphorylated by protein phosphatase 1 in its absence. When phosphorylated, DARPP-32 inhibits protein phosphatase 1, amplifying the effects of PKA-mediated phosphorylation of other substrates. DARPP-32 thus acts as a molecular switch that integrates signals from dopamine and other neurotransmitters and determines the overall response of the neuron.

\subsection{Eric Kandel and the Molecular Basis of Memory}

Kandel's contribution was the elucidation of the molecular mechanisms underlying learning and memory. He chose the marine invertebrate Aplysia californica---a sea slug with a relatively simple nervous system containing approximately 20,000 neurons---as his experimental model. Aplysia exhibits several forms of learning, including habituation (a decrease in response to repeated innocuous stimuli), sensitization (an enhanced response following a noxious stimulus), and classical conditioning.

Kandel showed that habituation of the gill-withdrawal reflex in Aplysia is caused by a decrease in the amount of neurotransmitter released from the sensory neuron at the synapse with the motor neuron. Sensitization, in contrast, is caused by the activation of a modulatory interneuron that releases serotonin onto the sensory neuron. Serotonin activates adenylyl cyclase in the sensory neuron, increasing cAMP levels, which activates PKA. PKA phosphorylates potassium channels, reducing potassium currents and broadening the action potential, which increases calcium influx and neurotransmitter release. This enhances the synaptic connection between the sensory and motor neurons.

Kandel further showed that long-term memory---which requires new protein synthesis---involves the activation of a transcription factor called CREB (cAMP response element-binding protein). CREB, activated by PKA, binds to specific DNA sequences (CRE elements) in the promoters of target genes and activates the transcription of genes required for the growth of new synaptic connections. Long-term memory is thus accompanied by structural changes at the synapse---the growth of new synaptic boutons and the strengthening of existing connections.

\section{Impact on Modern Medicine}

The discoveries of Carlsson, Greengard, and Kandel have had a significant impact on neuroscience and medicine.

\subsection{Neurodegenerative Disease}

Carlsson's identification of dopamine deficiency as the cause of Parkinson's disease led to the development of L-DOPA and dopamine agonist therapies that have improved the lives of millions of patients. Current research on Parkinson's disease, including deep brain stimulation, gene therapy, and cell-based therapies aimed at replacing lost dopamine neurons, is directly informed by Carlsson's foundational work.

\subsection{Psychiatric Disorders}

The dopamine hypothesis of schizophrenia, proposed by Carlsson, guided the development of first-generation and second-generation antipsychotic drugs. The second messenger cascades elucidated by Greengard have informed the development of new antidepressant drugs, including the selective serotonin reuptake inhibitors (SSRIs) and the phosphodiesterase inhibitors that are being investigated as novel antidepressants.

\subsection{Learning and Memory Disorders}

Kandel's work on the molecular mechanisms of learning and memory has implications for the treatment of Alzheimer's disease, age-related memory impairment, and other cognitive disorders. Understanding the signaling pathways that underlie synaptic plasticity provides targets for therapeutic intervention. Phosphodiesterase inhibitors that enhance cAMP signaling are being tested as potential treatments for cognitive impairment.

\subsection{Drug Addiction}

The signaling pathways discovered by Greengard and Kandel are also implicated in drug addiction. Drugs of abuse---cocaine, amphetamines, opioids, and nicotine---all affect dopamine and other neurotransmitter signaling pathways, and the synaptic plasticity mechanisms that underlie learning are hijacked by addictive drugs to create lasting changes in brain circuitry. Understanding these mechanisms is guiding the development of new treatments for addiction.

\section{Legacy}

Arvid Carlsson continued his research at the University of Gothenburg until his retirement. He received numerous honors, including the Wolf Prize in Medicine in 1979 and the Shaw Prize in Life Science and Medicine in 2004. Carlsson passed away on June 29, 2018, at the age of 95.

Paul Greengard spent his entire career at Rockefeller University, where he became Vincent Astor Professor and director of the Laboratory of Molecular and Cellular Neuroscience. He received the Albert Lasker Basic Medical Research Award in 1994. Greengard passed away on April 13, 2019, at the age of 96. He established the Greengard Prize, an annual award supporting women in science.

Eric Kandel continued his work at Columbia University, where he became university professor and director of the Kavli Institute for Brain Science. He has authored several influential books on the biology of memory, including \textit{In Search of Memory} (2006), which won the Aventis Prize for Science Books. Kandel continues to study the molecular mechanisms of memory formation and the cognitive neuroscience of consciousness.

Together, Carlsson, Greengard, and Kandel revealed the chemical logic of nervous system signaling---from the identification of neurotransmitters and their receptors to the intracellular cascades that convert chemical signals into lasting changes in synaptic strength. Their work has fundamentally shaped our understanding of how the brain processes information, forms memories, and is disrupted in neurological and psychiatric disease.


\section{Modern Developments}

Carlsson, Greengard, and Kandel's signal transduction work has led to modern neuroscience. Understanding of dopamine signaling has led to treatments for Parkinson's disease (levodopa) and schizophrenia (antipsychotics). Understanding of synaptic plasticity has informed development of cognitive enhancers. Understanding of CREB signaling has led to research on memory enhancement. Understanding of serotonin signaling has enabled development of antidepressants (SSRIs, SNRIs).


\section{Bibliography}

\begin{thebibliography}{9}
\bibitem{nobel-2000}
Nobel Prize Outreach, ``The Nobel Prize in Physiology or Medicine 2000,''\emph{NobelPrize.org}.\\
\url{https://www.nobelprize.org/prizes/medicine/2000/summary/}

\bibitem{lecture-2000}
Arvid Carlsson, ``Nobel Lecture,''\emph{NobelPrize.org}. Winner-authored primary source; downloadable from the official Nobel Prize archive.\\
\url{https://www.nobelprize.org/prizes/medicine/2000/carlsson/lecture/}
\end{thebibliography}
